\chapter[Sermon 4. Of the Narration of This Book, Where Also Is Discoursed of the Place and Time, and of the Author of This Revelation.]{Of the Narration of This Book, Where Also Is Discoursed of the Place and Time, and of the Author of This Revelation.}
\label{sermon:004}
\markright{Sermon 4. Of the Narration of This Book, Where Also Is ...}

\index{Patmos}I, John, your brother and companion in tribulation, and in the kingdom and patience which are in Jesus Christ, was on the island of Patmos for the word of God and for the testimony of Jesus Christ. I was in the spirit on the Lord’s Day and heard behind me a great voice, as if of a trumpet, saying, “I am Alpha and Omega, the first and the last.” And I saw another vision, and it was commanded that I write in a book and send it to the congregations which are in Asia, to Ephesus, and to Smyrna, and to Pergamum, and to Thyatira, and to Sardis, and to Philadelphia, and to Laodicea.

Note 4.1: The final section of the first part shows us a brief account wherein the Apostle John declares the time and place of this Revelation, and by whose authority he sent it to the seven churches in Asia.

And again now for the third time the name of John is repeated. He saw undoubtedly that there would be some who, to take away the use and benefit of this book, would doubt the author. Against whom he repeats and reiterates his name so often, lest we should doubt and lose the great value of so worthy a book.

He adds certain things to his name, which instruct concerning the state of the Apostle, and certain profitable matters. First, he calls himself a brother, namely of those churches, and of all ours. As where I have admonished you that in the seventh number are comprised all churches of all times throughout the whole world. We are all, so many believe, the children of one heavenly father. And therefore all spiritual brethren in Christ, coheirs with Christ, and heirs of God. Which thing Paul taught after Christ.

And seeing our dignity is so great, let us be ashamed of our misdeeds, lest our memory be put out of this most noble and celestial family. It is a shame that a brother of Christ, of John, and all the Apostles should degenerate. And so forth.

But why have they not so instantly urged this brotherhood, as the Monks have promoted their forged fraternities, the Rosaries of the Virgin Mary and of Saints? Because that was free, and cost nothing. But the Monks sell theirs dearly. They are therefore deceivers and seducers.

After he calls himself a partaker in affliction, or oppression and persecution, as he who was even now banished by the Emperor Domitian and lived in exile. And he joins together and does not separate himself in the evil common to all the faithful brethren. And truly it is one and the same persecution that vexed the Apostles and torments us at this day. Let us therefore rejoice that we have the Apostles and all the Martyrs of Christ as fellows in our trouble and affliction, that we are broken and bruised with the heavy burden of evils. Let us therefore be patient and long-suffering. For it is not enough to be afflicted and vexed with all kinds of evils (for many without any fruit or praise at all endure most grievous pains). But it becomes us also to be patient in adversity. Therefore John at this present joins with all in patience. For the Lord said in the Gospel, “In your patience you shall possess your souls.”

After he adds unto tribulation and patience a kingdom, and that an heavenly not a terrestrial kingdom. And he brings in the kingdom for the comfort of the patient people. For also the Apostle Paul said, a certain and sure saying. For if we die with Christ, we shall live also with him: If we suffer, we shall reign with him. \&c. Let us always here with comfort our selves in adversity. For we are thrust down, that we might once be exalted again, 2 Corinthians 4. And all these things are concluded in Christ Jesus, by whom we are both the children and brethren of God, and suffer many things patiently, and are made partakers of his kingdom. For even for these things must we thank him and his mercies, and not our own desert.

Let us note here also what and how great the humility of the greatest and worthy Apostle of God was, considering his state: It was not pleasant, but difficult, yet in patience joyful through Christ. But where are those now who glory in the name of Apostles? Who, meanwhile, swelling with pride, are addicted to filthy pleasures? I warn you, that we flee from them as from apostates.

And now he shows the place where this divine revelation was made to him, where also he was commanded by God to write the same. The place was the Isle of Patmos. It is accounted among the islands called Sporades by Pliny in the fourth book and twelfth chapter. It lay opposite Asia and the city of Ephesus, and was in the sight both of Europe and Africa, so that it seemed to be as it were a middle seat, or holy chair, out of which Christ preached through John from heaven to the whole world. And indeed the councils of God are wonderful, and his goodness is unspeakable, which reveals so great mysteries, as it were in the Roman prison or Babylonian captivity, to his faithful.

\index{Tertullian}\index{Eusebius of Caesarea}\index{Rome}Neither does he hide the cause of his coming into the island, saying, “I was there for the word of God and the testimony of Jesus Christ.” The word of God is the very essence of God, called by John with a unique expression the word or sermon of God, as appears in the first chapter of John. The testimony of Jesus Christ is the Gospel itself, which Jesus testified, and which his disciples have testified concerning Jesus. Therefore, for the confession and preaching of Jesus Christ and his wholesome Gospel (for he also explained how he became a partaker of affliction), John was apprehended in Asia and led by soldiers to Rome, that he might plead his cause before Emperor Domitian, who, by his cruel nature, condemned the innocent. And he was put into a cauldron of hot boiling oil. When he escaped unharmed from this, he was carried to Patmos. He answered no other matter before the emperor than Paul did twenty-seven years before, during Nero’s reign. This occurred in the thirteenth or fourteenth year of Domitian’s reign and twenty-four years after the destruction of the city of Jerusalem and sixty-six years after the birth of our Lord. Domitian, who sought to be called a god, was slain by his own men after many murders and cruel acts, and died himself a shameful death in the fifteenth year of his reign. The authorities for this are Suetonius in the life of Domitian, Tertullian in \textit{Prescription Against Heretics}, Eusebius in his chronicles and in the third book of the \textit{Ecclesiastical History}, chapters seventeen and eighteen. And to this is added the common consent of all writers.

\index{Resurrection}Moreover, he notes the time also, in which these mysteries began to be revealed to him, [illegible], in that solemn day of the Lord, namely Sunday. For so the ancient fathers have called one of the Sabbaths, that is to say, the first day of the week, wherein Christ rose again from the dead, as recorded in Matthew 28 and Mark 16. And this day the churches have chosen for themselves in place of the Sabbath day, as holy in remembrance of the Lord’s resurrection, wherein they might keep their sacred and solemn assemblies. For that this day was observed and consecrated for assemblies in the congregation of Corinth appears manifestly in the 16th chapter of the first Epistle to the Corinthians, where the Apostle commands them to lay apart their collections on one of the Sabbaths. The same day also the faithful did celebrate their service with Paul, as recorded in Acts 20. Where Sozomenus reports in the 8th chapter of the first book of the \textit{Ecclesiastical History}, that great Constantine made certain holy days, and even the Lord’s day, which the Heathens call Sunday; it is to be understood that he renewed rather the custom of the Apostles and the catholic church, than that he newly instituted it. And freely the churches have received that day; for we read not that it was anywhere commanded. And the congregations saw how it was altogether necessary that there should be a certain time in which the saints should meet and come together. They chose therefore the day of the resurrection, neither did they maliciously contend among themselves for these things, as the histories testify was done in the church afterward. And at this day, verily, the superstitious holy days being abrogated, it is better to observe certain and moderate days, and to keep peace and quietness in the church.

But where this Apostle knew that the faithful on that day served God in all assemblies, where he could not be present in body, in spirit and contemplation he was with them. And as he was thus in the spirit and contemplation of divine matters, and in holy prayers, he heard a voice, whereof I will speak hereafter. But here we are presently taught what is the religion of the Sunday, and how it is fitting to serve it. Finally, worldly men are reproved who pollute and break it with profane works and affairs. David, in his time, when he suffered persecution from Saul, lamented chiefly that he might not come to the Lord’s tabernacle. Our people, however, count it a great felicity never to enter into the fellowship of the Saints. And to abuse the Sunday, in gaming, drinking, dancing, and worldly business.

These things declared in this manner, he comes now to the Revelation, setting forth before the clear commandment of God, by which he was commanded both to write the things revealed and also to send them to the seven churches of Asia. A chief characteristic of the Revelation’s manner and majesty is that he heard a voice, and one notable as the sound of a trumpet. For so we read it was done in the giving of the law at Mount Sinai. Now it is declared who that voice was, and who the author of the Revelation. Truly, it was the eternal God, who calls himself Alpha and Omega, that is, the beginning and the end. Or as it is said in Essay, first and last, as we shall see elsewhere.

\index{Apocalypse (Book of)}Now follows the commandment which has several parts. For first, the Lord commands John to see and to write such things as he saw, that is to wit, the Apocalypse. And he should write neither on a tablet nor on the wall, but in a book: Verily for the edifying and profit of the church present, and of all posterity. After this, he also commanded to send those writings to seven congregations, and verily to all the churches of the whole world for all times and ages. Therefore, all these things belong to the provision of congregations, and that of all that be, have been, or shall be.

Here we learn how great is the authority of the scriptures. It was not written nor compiled in books, but by God’s commandment. There are notable testimonies of the book of Moses, in the thirty-fourth chapter of Exodus and the thirty-first chapter of Deuteronomy. And to say nothing of the rest of the prophets, was not Jeremiah commanded to write his sermons again, which King Jehoiakim had cut in pieces and burnt?

Doubts Saint Peter bears manifest witness that the prophets received the mysteries of God to no other end than that they should reveal them to us; which indeed might only be done by the scriptures. Now John is most clearly commanded to write. What will we say that he is also commanded to send his writings to the congregations? Whereof again we gather that God wills right well to the congregations, and even to every one of us. Let us beware and take heed that we do not put from us unworthily so great benefits of God, to whom be praise and glory.
